In this paper, we introduced a novel neighborhood search operator tailored for the Just-In-Time Job Shop Scheduling (JITJSS) problem with earliness and tardiness penalties, designed to exploit critical-path structure and resolve scheduling bottlenecks with minimal disruption to adjacent operations.
To eliminate the subjective trial-and-error often associated with metaheuristic design and objectively assess the merit of our operator, we embedded it within an Automatic Algorithm Design (AAD) framework.
By formalizing the algorithmic design space of Iterated Local Search (ILS) via a Backus--Naur Form (BNF) grammar and configuring it using \texttt{irace}, the automated design process provided independent, data-driven confirmation of the operator's value: the configured elite pipelines systematically selected and prioritized our proposed neighborhood, demonstrating its central role in achieving competitive search performance.

Rigorous computational experiments across the 72 benchmark instances of \citet{baptiste2008lagrangian}, conducted under conservative runtime budgets calibrated through single-thread PassMark benchmarking, established the effectiveness of the resulting algorithm, $C_1$.
The configured solver matched or improved the Best-Known Solutions (BKS) in 72.2\% of instances, exhibiting consistent performance across problem sizes (70.8\% for $n=10$, 75.0\% for $n=15$, and 70.8\% for $n=20$).
Pairwise one-tailed Wilcoxon signed-rank tests confirmed that $C_1$ significantly outperforms all literature baselines covering the complete benchmark set ($p < 0.05$), while maintaining statistical comparability ($p = 0.09912$ and $p = 0.06681$) with recent adaptive and reinforcement learning methods on small and medium instances.
Importantly, our approach maintained practical runtime scalability on large 20-job instances where complex matheuristics and neural models encounter severe computational bottlenecks, converging to final best solutions in an average of 3.7, 20.7, and 61.4 seconds across problem scales.
Beyond the standalone performance of $C_1$, the collective search across elite configurations established tighter upper bounds for 25 benchmark instances, updating historical records for future benchmarking studies.

These results confirm that combining domain-specific neighborhood engineering with automated algorithmic configuration yields robust, scalable solvers that match or surpass manually tuned state-of-the-art methods.
Natural extensions of this work include integrating the proposed operator into multi-objective settings with sequence-dependent setup times, evaluating its responsiveness in dynamic reactive rescheduling environments, and investigating automated selection mechanisms to exploit complementary dynamics among elite configurations.
